\chapter{Giới thiệu}
\label{chap:intro}

\section{Bài toán và động lực}

Đái tháo đường (ĐTĐ, bệnh tiểu đường) là tình trạng glucose trong máu tăng cao kéo dài do cơ thể thiếu insulin hoặc kháng insulin. Phần lớn ca bệnh là ĐTĐ típ 2, loại tiến triển âm thầm nhiều năm và có thể làm chậm hoặc phòng ngừa nếu phát hiện sớm từ giai đoạn tiền ĐTĐ \cite{ada2024,knowler2002dpp}. Ở Việt Nam, Hướng dẫn chẩn đoán và điều trị ĐTĐ típ 2 của Bộ Y tế \cite{byt5481} khuyến cáo sàng lọc chủ động cho người có yếu tố nguy cơ, nhưng người dân thường chỉ có trong tay vài con số rời rạc (một lần thử đường huyết, cân nặng, huyết áp) và không biết chúng nói lên điều gì.

Một công cụ hỗ trợ sàng lọc tốt phải làm được bốn việc cùng lúc: \emph{dự đoán} nguy cơ từ dữ liệu sức khỏe; \emph{giải thích} dựa trên tri thức y khoa có nguồn; \emph{gợi ý} bước can thiệp cụ thể; và \emph{cảnh báo} ngay khi có dấu hiệu nguy hiểm. Mỗi việc cần một kỹ thuật khác nhau: học sâu cho dự đoán, cơ sở tri thức và RAG cho giải thích, học tăng cường cho ra quyết định tuần tự, và luật cố định cho an toàn. Gộp tất cả vào một chương trình khiến hệ thống khó kiểm tra và khó tin cậy. Vì vậy project xây dựng một \textbf{hệ đa agent}: mỗi agent giữ một chuyên môn, các agent trao đổi bằng thông điệp có cấu trúc, và một hội đồng kiểm tra chéo kết luận trước khi tư vấn cho người dùng.

\section{Mục tiêu và phạm vi}

Project hiện thực hóa sơ đồ mà giảng viên đưa ra trên lớp (\figref{fig:board}, Chương~\ref{chap:arch}) thành một hệ thống chạy được đầu–cuối:
\begin{itemize}
  \item \textbf{DB → Deep Learning → dự đoán, phân loại}: huấn luyện mô hình học máy và mạng nơ-ron trên hai bộ dữ liệu có sẵn của CDC, tổng cộng \textbf{3.193.465 người}: NHANES 2011–2018 (20.646 người, nhãn xác nhận bằng xét nghiệm) làm mô hình chính, và BRFSS 2011–2023 (\textbf{3.172.819 người}, bảy năm khảo sát) làm mô hình quy mô lớn; ước lượng xác suất ĐTĐ rồi phân tầng nguy cơ thành bốn mức.
  \item \textbf{Web, sách → Neo4j KB → RAG}: nạp hướng dẫn ADA 2024 và QĐ 5481/QĐ-BYT vào cơ sở tri thức; đồ thị tri thức lưu trên \textbf{Neo4j}, đoạn văn được tìm kiếm lai trên \textbf{Qdrant} (vector) và \textbf{SQLite FTS5} (toàn văn).
  \item \textbf{Reinforcement learning + rule → Agent}: agent gợi ý can thiệp học bằng Q-learning nhưng bị ràng buộc bởi luật an toàn và kết luận của hội đồng.
  \item \textbf{Tư vấn, trợ giúp}: hội đồng tranh biện liệt kê bằng chứng đồng thuận (+) và mâu thuẫn (−), rồi agent phản hồi viết câu trả lời có trích dẫn.
  \item \textbf{Code + Demo}: chatbot web kèm \emph{Phòng hội chẩn} trực quan hóa từng thông điệp giữa các agent; mã nguồn có kiểm thử tự động.
\end{itemize}

\textbf{Phạm vi.} Đây là sản phẩm học tập. Hệ thống phục vụ \emph{sàng lọc} chứ không \emph{chẩn đoán}, không đưa liều thuốc cụ thể, và mọi câu trả lời đều nhắc người dùng đi khám bác sĩ. Mọi con số trong báo cáo được đo lại từ mã nguồn của project.

\section{Đóng góp và cấu trúc báo cáo}

Các đóng góp chính:
\begin{itemize}
  \item Một hệ \textbf{chín agent} giao tiếp qua \texttt{MessageBus} theo tinh thần FIPA-ACL, có audit trail đầy đủ và cô lập lỗi.
  \item Khối DB gồm \textbf{hơn 3,19 triệu hồ sơ} thật của CDC, gấp hơn 12 lần bản BRFSS 2015 đã làm sạch (253.680 dòng) thường được dùng cho bài toán này. Mô hình huấn luyện trên dữ liệu quy mô lớn đóng vai \emph{ý kiến thứ hai}, và hội đồng chất vấn khi nó ngược với mô hình chính.
  \item Hai nhánh song song đúng như sơ đồ: \emph{DB → ML/DL} và \emph{KB (Neo4j) → RAG}, gặp nhau ở một hội đồng có \textbf{tranh biện nhiều vòng}: agent phản biện chất vấn agent chuyên môn khi chúng mâu thuẫn.
  \item Agent \textbf{RL + luật}: chính sách Q-learning chỉ được hành động trong khuôn khổ do tri thức tường minh cho phép.
  \item \textbf{Phòng hội chẩn}: giao diện phát lại luồng thông điệp, cho người xem thấy agent nào đang nói gì với agent nào.
  \item Đánh giá trên 634.564 hồ sơ BRFSS kiểm tra và đánh giá đầu–cuối trên 400 ca NHANES chưa từng thấy; đo lại truy xuất RAG trên bộ lưu trữ mới.
\end{itemize}

Báo cáo đi theo đúng dàn ý trên bảng. Chương~\ref{chap:mas} trình bày cơ sở về agent và hệ đa agent. Chương~\ref{chap:arch} là kiến trúc tổng quan. Chương~\ref{chap:comp} mô tả từng thành phần: Deep Learning, KB và RAG, RL kết hợp luật, hội đồng và tầng tiếp nhận. Chương~\ref{chap:code} trình bày mã nguồn, demo, kiểm thử và đánh giá. Cuối cùng là kết luận.
